% from ega4-2-fr.tex
\oldpage[IV-2]{159}

où les $L_i$ sont des $A$-modules libres de type fini ($A$ étant
noethérien), d'où une suite exacte
\[
0\longrightarrow R\longrightarrow L_{n-1}\longrightarrow L_{n-2}
\longrightarrow\cdots\longrightarrow L_0\longrightarrow M
\longrightarrow0
\]
où $R$ est un $A$-module de type fini; on en déduit une suite exacte
\[
0\longrightarrow R_x\longrightarrow(L_{n-1})_x\longrightarrow\cdots
\longrightarrow(L_0)_x\longrightarrow M_x\longrightarrow0
\]
où les $(L_i)_x$ sont des $A_x$-modules libres de type fini; comme par
hypothèse $\operatorname{dim.proj}(M_x)=n$, cela entraîne que $R_x$ est un
$A_x$-module projectif de type fini (M, VI, 2.1) et par suite un
$A_x$-module libre de type fini (0$_{\rm III}$, 10.1.3). On en conclut qu'il
existe un voisinage ouvert $U$ de $x$ dans $X$ tel que, pour tout
$x'\in U$, $R_{x'}$ soit un $A_{x'}$-module libre (0$_{\rm I}$, 5.2.7),
donc $M_{x'}$ admet une résolution projective de longueur $n$, autrement
dit $\operatorname{dim.proj}(M_{x'})\leq n$, ce qui prouve le lemme.
\end{proof}

\begin{proposition}[6.11.2]
\label{IV.2.s6.11.2-fr}
Soit $X$ un préschéma localement noethérien et localement immersible dans
un schéma régulier (5.11.1); soit $\mathcal F$ un $\mathcal O_X$-Module
cohérent. Alors:
\begin{enumerate}
\item[(i)] (M. Auslander) La fonction
$x\mapsto\operatorname{coprof}(\mathcal F_x)$ est semi-continue
supérieurement dans $X$ (autrement dit, pour tout entier $n$, l'ensemble
$U_{C_n}(\mathcal F)$ des $x\in X$ tels que
$\operatorname{coprof}(\mathcal F_x)\leq n$ est ouvert).
\item[(ii)] Pour tout entier $n$, l'ensemble $U_{S_n}(\mathcal F)$ des
$x\in X$ tels que $\mathcal F$ possède la propriété $(S_n)$ au point $x$
est ouvert.
\end{enumerate}
\end{proposition}

\begin{proof}
\begin{enumerate}
\item[(i)] La question étant locale sur $X$, on peut, en vertu de
l'hypothèse, se borner au cas où $X$ est un sous-schéma fermé d'un schéma
affine régulier $Y$. Si $j:X\to Y$ est l'injection canonique, et
$\mathcal G=j_*(\mathcal F)$, on sait qu'on a alors, pour tout $x\in X$,
$\dim(\mathcal F_x)=\dim(\mathcal G_x)$ et
$\operatorname{prof}(\mathcal F_x)=\operatorname{prof}(\mathcal G_x)$,
donc $\operatorname{coprof}(\mathcal F_x)=
\operatorname{coprof}(\mathcal G_x)$, et comme
$\operatorname{coprof}(\mathcal G_y)=0$ pour $y\in Y-X$, on est ramené à
prouver la propriété pour $\mathcal G$; autrement dit, on peut se borner au
cas où $X$ est un schéma affine régulier. On sait alors (0, 17.3.4) que
l'on a
\[
\operatorname{prof}(\mathcal F_x)=\dim(\mathcal O_{X,x})-
\operatorname{dim.proj}(\mathcal F_x).
\]
D'autre part, si $S$ est l'unique sous-préschéma fermé de $X$ ayant pour
espace sous-jacent $\operatorname{Supp}(\mathcal F)$, on a
$\dim(\mathcal F_x)=\operatorname{codim}(\{x\},S)$ (5.1.12.1), et, en
vertu de (5.1.9)
\[
\dim(\mathcal F_x)=\dim(\mathcal O_{X,x})-\operatorname{codim}_x(S,X)
\]
puisque $\mathcal O_{X,x}$ est un anneau régulier, et a fortiori
biéquidimensionnel (0, 16.5.12). On peut donc écrire
\begin{equation}
\operatorname{coprof}(\mathcal F_x)=
\operatorname{dim.proj}(\mathcal F_x)-\operatorname{codim}_x(S,X)
\tag{6.11.2.1}\label{IV.2.s6.11.2.1-fr}
\end{equation}
et la proposition résulte alors de (6.11.1) et de (0, 14.2.6).

